\documentclass{article}

\usepackage[preprint]{neurips_2025}

\usepackage[utf8]{inputenc}
\usepackage[T1]{fontenc}
\usepackage{hyperref}
\usepackage{url}
\usepackage{booktabs}
\usepackage{amsfonts}
\usepackage{amsmath}
\usepackage{nicefrac}
\usepackage{microtype}
\usepackage{xcolor}
\usepackage{enumitem}

\title{RL-based Dynamic Scheduling for LLM Serving:\\Improving Llumnix with Reinforcement Learning}

\author{
  Yuxuan Bai \\
  University of Virginia \\
  \texttt{yuxuan.bai@virginia.edu} \\
  \And
  Jiaqi Liu \\
  University of Virginia \\
  \texttt{jiaqi.liu@virginia.edu} \\
  \And
  Peng Xia \\
  University of Virginia \\
  \texttt{peng.xia@virginia.edu} \\
}

\begin{document}

\maketitle

\vspace{-1em}
\section{Introduction}
\vspace{-0.3em}

Serving large language models (LLMs) at scale presents unique scheduling challenges.
Unlike fixed-size inference tasks, LLM requests exhibit \emph{unpredictable output lengths}, \emph{dynamically growing KV caches}, and \emph{heterogeneous latency requirements}, making static scheduling fundamentally inadequate.
When GPU memory is exhausted, the system must \emph{preempt} running requests, causing severe tail-latency degradation (up to 3.8$\times$ P99 increase). In multi-GPU clusters, load imbalance and memory fragmentation further compound these issues.

Llumnix~\cite{sun2024llumnix} addresses these challenges by introducing \emph{runtime rescheduling}---live migration of running requests with their KV caches across GPU instances via pipelined transfers with near-zero downtime.
However, Llumnix relies on \emph{hand-crafted heuristics} (a ``virtual usage'' mechanism) to decide when and where to migrate. These heuristics are static, myopic, and use fixed tradeoff parameters, limiting their effectiveness under dynamic workloads.

We propose replacing Llumnix's heuristic policy with a \textbf{reinforcement learning (RL)-based policy} that learns scheduling decisions from data. Our RL agent observes cluster state and outputs virtual usage values that drive Llumnix's existing infrastructure. This approach is \emph{adaptive} (learns from workload patterns), \emph{data-driven} (eliminates manual rule crafting), and \emph{globally optimizing} (maximizes long-term cumulative reward). Critically, our design requires \textbf{minimal modification}: only the \texttt{CalcVirtualUsage()} function is replaced, while migration, dispatch, and the vLLM backend remain unchanged.

\vspace{-0.5em}
\section{Background}
\vspace{-0.3em}

\paragraph{LLM serving and KV cache management.}
Auto-regressive LLM inference generates tokens sequentially, with each token attending to all previous tokens via key-value (KV) caches in GPU memory. The KV cache grows linearly with output length, which is unknown at dispatch time. Systems like vLLM~\cite{kwon2023vllm} use paged attention for efficient memory management, but when aggregate usage exceeds GPU capacity, requests must be preempted---swapped to CPU or recomputed---causing severe latency spikes.

\paragraph{Llumnix and its virtual usage heuristic.}
Llumnix~\cite{sun2024llumnix} (OSDI 2024) is the first system to support live migration of LLM requests across GPUs. It introduces a \emph{virtual usage} abstraction that may differ from physical usage to unify scheduling objectives. A ``freeness'' score $F_i = (M - \sum \text{Virtual}_j) / B$ determines request placement. Four hand-crafted rules govern virtual usage: (1)~normal requests use physical usage; (2)~queuing requests inflate to demand level; (3)~high-priority instances add headroom; (4)~terminating instances set usage to $\infty$. While effective, these heuristics are \emph{static} (manually crafted, cannot adapt to new workloads), \emph{myopic} (no predictive lookahead, greedy local optima), and have \emph{fixed tradeoffs} (hardcoded balance among competing objectives).

\paragraph{RL for systems scheduling.}
RL has demonstrated success in cluster resource management~\cite{mao2016resource}, video streaming~\cite{mao2017neural}, and query optimization~\cite{marcus2018towards}, outperforming hand-tuned heuristics by learning policies that account for long-term consequences. Our work applies this paradigm to LLM serving, where dynamic KV cache growth and live migration create a well-suited MDP structure.

\vspace{-0.5em}
\section{Project Tasks}
\vspace{-0.3em}

We decompose the project into five tasks with clear ownership.

\paragraph{Task 1: MDP formulation and environment design \textnormal{(Yuxuan Bai).}}
Define the scheduling MDP. The state $s_t$ captures per-instance features (memory usage, queue length), per-request features (token counts, priority class, waiting time), and cluster metrics (total load, fragmentation). The action $a_t$ outputs virtual usage values as a drop-in replacement for \texttt{CalcVirtualUsage()}. The reward $r_t$ is a weighted combination minimizing P99 latency, preemptions, and fragmentation while maximizing SLO compliance and cost efficiency. This task includes building a simulation environment interfacing with Llumnix's scheduling loop.

\paragraph{Task 2: RL agent implementation and training \textnormal{(Jiaqi Liu).}}
Implement the RL policy network and training pipeline using PPO~\cite{schulman2017proximal} for its stability in continuous action spaces. The policy uses a multi-head architecture outputting virtual usage values per instance. Training proceeds in two phases: offline pre-training from heuristic scheduling traces, then online fine-tuning in the simulated environment. We also implement baseline comparisons (original heuristic, DQN with discretized actions).

\paragraph{Task 3: Llumnix integration \textnormal{(Peng Xia).}}
Integrate the RL policy into Llumnix, replacing \texttt{CalcVirtualUsage()} with policy inference. This includes: (a)~a state collector extracting real-time cluster features; (b)~low-latency policy deployment ($< 1$ms per decision); (c)~correctness validation by comparing traces between heuristic and RL modes under identical workloads.

\paragraph{Task 4: Experimental evaluation \textnormal{(All members).}}
Conduct five experiments on 16$\times$ A10 GPUs (24GB) with vLLM backend, comparing against Round-Robin, INFaaS++, and Llumnix-VU baselines.
\emph{Exp~1}: Baseline serving on ShareGPT/BurstGPT traces (mean \& P99 latency).
\emph{Exp~2}: Priority workloads with 10\% high-priority requests (SLO isolation).
\emph{Exp~3}: Bursty traffic with high-CV Gamma arrivals (adaptation under spikes).
\emph{Exp~4}: Long-context workloads stressing memory fragmentation.
\emph{Exp~5}: Generalization---train on ShareGPT, test on BurstGPT.

\paragraph{Task 5: Analysis and report \textnormal{(All members).}}
Analyze results including ablation studies on reward weights and state features. Prepare the final report with findings, limitations, and future directions (multi-objective RL, direct migration actions).

\begin{thebibliography}{6}
\small

\bibitem[Sun et al.(2024)]{sun2024llumnix}
B.~Sun, Z.~Huang, H.~Zhao, W.~Xiao, X.~Zhang, Y.~Li, and W.~Lin.
Llumnix: Dynamic scheduling for large language model serving.
In \emph{OSDI}, 2024.

\bibitem[Kwon et al.(2023)]{kwon2023vllm}
W.~Kwon, Z.~Li, S.~Zhuang, Y.~Sheng, L.~Zheng, C.~H.~Yu, J.~E.~Gonzalez, H.~Zhang, and I.~Stoica.
Efficient memory management for large language model serving with {PagedAttention}.
In \emph{SOSP}, 2023.

\bibitem[Mao et al.(2016)]{mao2016resource}
H.~Mao, M.~Alizadeh, I.~Menache, and S.~Kandula.
Resource management with deep reinforcement learning.
In \emph{HotNets}, 2016.

\bibitem[Mao et al.(2017)]{mao2017neural}
H.~Mao, R.~Netravali, and M.~Alizadeh.
Neural adaptive video streaming with {Pensieve}.
In \emph{SIGCOMM}, 2017.

\bibitem[Marcus et al.(2018)]{marcus2018towards}
R.~Marcus, P.~Negi, H.~Mao, C.~Zhang, M.~Alizadeh, T.~Kraska, and N.~Tatbul.
Towards learned infrastructure.
\emph{arXiv:1812.01772}, 2018.

\bibitem[Schulman et al.(2017)]{schulman2017proximal}
J.~Schulman, F.~Wolski, P.~Dhariwal, A.~Radford, and O.~Klimov.
Proximal policy optimization algorithms.
\emph{arXiv:1707.06347}, 2017.

\end{thebibliography}

\end{document}
